\documentclass[11pt,a4paper]{article}
\usepackage[margin=1in]{geometry}
\usepackage{amsmath,amssymb,amsthm}
\usepackage{graphicx}
\usepackage{hyperref}
\usepackage{booktabs}
\usepackage{float}
\usepackage{enumitem}
\usepackage{xcolor}
\usepackage{svg}
\svgsetup{inkscapelatex=false,inkscapeformat=pdf,inkscapepath=svgdir}

\hypersetup{colorlinks=true,linkcolor=blue,urlcolor=cyan,citecolor=red}

\newtheorem{theorem}{Theorem}[section]
\newtheorem{proposition}[theorem]{Proposition}
\newtheorem{corollary}[theorem]{Corollary}
\newtheorem{remark}[theorem]{Remark}
\newtheorem{definition}[theorem]{Definition}

\newcommand{\ket}[1]{|#1\rangle}
\newcommand{\bra}[1]{\langle#1|}
\newcommand{\braket}[2]{\langle#1|#2\rangle}
\newcommand{\SSV}{\mathrm{SSV}}

\title{\textbf{Quantum Mechanics in Conscious Point Physics:\\
Entanglement and Bell Inequality Violation\\
from Non-Separable Nexus States}\\[6pt]
{\large QM Series --- Paper 4 (revised)}}

\author{Thomas Lee Abshier, ND\thanks{Hyperphysics Institute;
\texttt{drthomas007@protonmail.com};
\texttt{https://hyperphysics.com}} \\
\small{CPP-2040c --- companion to CPP-2040a (overview),
CPP-2040b (Schr\"{o}dinger), CPP-041 (QM series)}}

\date{22 March 2026 --- Version~2}

\begin{document}
\maketitle

\begin{abstract}
We derive quantum entanglement and Bell inequality violation in
Conscious Point Physics (CPP) from first principles.  The key result
is a \textbf{theorem} establishing why CPP is not a local hidden
variable (LHV) theory even though it is a deterministic, discrete
lattice model: the DI-bit representation of the singlet state is
\emph{non-separable} --- it cannot be factored into independent
states for particle~A and particle~B.  The Nexus maintains this
non-separability as particles separate, providing the non-local
resource required for Bell violation.  We give an explicit derivation
of the singlet correlation function $E(\hat{a},\hat{b}) = -\cos\theta$
from the Born rule (CPP-2040b, companion~C3~\citep{abshier2026born})
applied to the non-separable DI-bit state, and show this yields
$|\mathrm{CHSH}| = 2\sqrt{2}$ at optimal settings.  We prove the
no-signaling theorem holds in CPP: neither party's marginal probability
depends on the other's measurement axis.  We identify clearly what the
\emph{old} argument (shared bit pool with pre-established correlations)
got wrong, why it is equivalent to the EPR argument that Bell rules out,
and how the corrected argument differs.  The corrected CPP derivation
is consistent with all loophole-free Bell tests and makes one
falsifiable prediction: the correlation function acquires lattice
corrections of order $(\nu/\nu_P)^2$ at frequencies approaching the
Planck frequency $\nu_P = 1/t_P$, deviating from perfect quantum
correlations in a specific, calculable way.
\end{abstract}

\tableofcontents
\newpage

%===========================================================================
\section{Introduction and the Problem with the Previous Argument}
\label{sec:intro}
%===========================================================================

The previous version of this paper (QM Series~\#4, v1) attempted to
explain Bell inequality violation in CPP using the following argument:
EPR-like pairs share a \emph{fixed} DI-bit pool at emission; measurement
on one side depletes the pool; correlations arise from pre-established
phase relationships.

\textbf{This argument is wrong.}  It is structurally identical to the
Einstein--Podolsky--Rosen (EPR) local realism argument that Bell's
theorem explicitly rules out~\citep{bell1964}.  Bell proved that
\emph{any} theory where outcomes are determined by pre-existing
properties of the particles --- regardless of what those properties are
called --- cannot produce correlations stronger than $|\mathrm{CHSH}|
= 2$.  The quantum value is $2\sqrt{2} \approx 2.83$.

The error was treating the DI-bit phases as classical random variables
(hidden variables) that are set at emission and read out at measurement.
This gives the \emph{classical} correlation function
$E_{\rm cl}(\hat{a},\hat{b}) = 1 - 2\theta/\pi$, not the quantum value
$E_{\rm QM}(\hat{a},\hat{b}) = -\cos\theta$.  The classical model saturates
at $|\mathrm{CHSH}| = 2$; it never reaches $2\sqrt{2}$.

This paper provides the correct derivation.  The two key corrections are:
\begin{enumerate}
\item The DI-bit state of a singlet pair is \emph{non-separable}: it
  cannot be written as a product of independent states for the two
  particles.  The Nexus maintains this non-separability.
\item The probability rule is the Born rule ($P = |\text{projection}|^2$),
  not the classical rule ($P \propto |\text{projection}|$).  This is
  what produces $-\cos\theta$ rather than $1 - 2\theta/\pi$.
\end{enumerate}

%===========================================================================
\section{Bell's Theorem: What It Rules Out and What It Allows}
\label{sec:bell}
%===========================================================================

\subsection{The local hidden variable framework}

A local hidden variable (LHV) theory assumes:
\begin{itemize}
\item Each pair of particles shares a hidden variable $\lambda$
  distributed according to $\rho(\lambda)$.
\item Outcome at~A depends only on the local measurement setting $\hat{a}$
  and the local particle variable: $A = A(\hat{a}, \lambda)$.
\item Outcome at~B depends only on $\hat{b}$ and its local variable:
  $B = B(\hat{b}, \lambda)$.
\item The joint probability factorizes:
  $P(a,b|\hat{a},\hat{b},\lambda) = P_A(a|\hat{a},\lambda) \times
   P_B(b|\hat{b},\lambda)$.
\end{itemize}

\begin{theorem}[Bell--CHSH inequality~\citep{bell1964,clauser1969}]
\label{thm:bell}
For any LHV theory, the CHSH correlator satisfies
\begin{equation}
S = |E(\hat{a},\hat{b}) + E(\hat{a},\hat{b}') + E(\hat{a}',\hat{b})
    - E(\hat{a}',\hat{b}')| \;\leq\; 2.
\label{eq:chsh}
\end{equation}
\end{theorem}

\subsection{What escapes Bell's theorem}

Bell's theorem applies specifically to LHV models.  It does \emph{not}
apply to:
\begin{itemize}
\item \textbf{Non-separable states}: If the joint state of the two
  particles cannot be written as $|\psi_A\rangle \otimes |\psi_B\rangle$,
  the factorization assumption fails and Theorem~\ref{thm:bell} does
  not apply.
\item \textbf{Non-local constraints}: If the correlation mechanism
  involves a global constraint (not a local hidden variable), the
  factorization assumption fails.
\end{itemize}

Standard quantum mechanics uses non-separable states (entanglement).
CPP uses the Nexus to maintain non-separable DI-bit states.  Both
escape the Bell bound for the same structural reason.

%===========================================================================
\section{CPP Representation of the Singlet State}
\label{sec:singlet}
%===========================================================================

\subsection{Qubit structure of a spin-1/2 CP aggregate}

A spin-$1/2$ CP aggregate carries a two-component DI-bit state
$|\psi\rangle \in \mathbb{C}^2$ (a qubit).  The two basis states are
$|\!\uparrow\rangle$ (spin aligned with the quantization axis) and
$|\!\downarrow\rangle$ (anti-aligned).  The phase angle of the DI-bit
state determines the spin direction.

This qubit structure is the CPP representation of the standard
two-state quantum system.  It follows from the ZBW helical orbit of
the spin-orbital hDP (companion~C4~\citep{abshier2026zbw}): the orbit
axis defines the spin direction, and the complex phase of the orbit
is the DI-bit qubit phase.

\subsection{The singlet state as a non-separable joint DI-bit state}

When two spin-$1/2$ CP aggregates are produced with total spin zero
(a singlet), their joint DI-bit state is:
\begin{equation}
|\Psi^-\rangle \;=\;
\frac{1}{\sqrt{2}}\bigl(|\!\uparrow\rangle_A|\!\downarrow\rangle_B
                       -|\!\downarrow\rangle_A|\!\uparrow\rangle_B\bigr).
\label{eq:singlet}
\end{equation}

\begin{proposition}[Non-separability of the singlet]
\label{prop:nonsep}
The state~\eqref{eq:singlet} cannot be written as $|\phi_A\rangle \otimes
|\phi_B\rangle$ for any single-particle states $|\phi_A\rangle$ and
$|\phi_B\rangle$.
\end{proposition}

\begin{proof}
Suppose $|\Psi^-\rangle = |\phi_A\rangle \otimes |\phi_B\rangle$ with
$|\phi_A\rangle = \alpha|\!\uparrow\rangle + \beta|\!\downarrow\rangle$
and $|\phi_B\rangle = \gamma|\!\uparrow\rangle + \delta|\!\downarrow\rangle$.
Expanding: $|\phi_A\rangle \otimes |\phi_B\rangle =
\alpha\gamma|\uparrow\uparrow\rangle + \alpha\delta|\uparrow\downarrow\rangle
+ \beta\gamma|\downarrow\uparrow\rangle + \beta\delta|\downarrow\downarrow\rangle$.
Matching to~\eqref{eq:singlet}: $\alpha\gamma = 0$, $\alpha\delta = 1/\sqrt{2}$,
$\beta\gamma = -1/\sqrt{2}$, $\beta\delta = 0$.
From the first and last equations, either $\alpha = 0$ or $\gamma = 0$, and
either $\beta = 0$ or $\delta = 0$.  But then either $\alpha\delta = 0$ or
$\beta\gamma = 0$, contradicting $|\alpha\delta| = |\beta\gamma| = 1/\sqrt{2}$.
\end{proof}

\begin{remark}[Physical interpretation]
The non-separability means the two CP aggregates do not carry
independent spin states.  They share a joint DI-bit state.  There
is no way to assign a definite spin direction to particle~A alone
or to particle~B alone --- the spin direction of each is only defined
relative to the other, through the joint state.
\end{remark}

\subsection{The Nexus maintains non-separability}

As the two CP aggregates travel apart after emission, the joint DI-bit
state~\eqref{eq:singlet} is maintained by the Nexus.  At each Absolute
Moment tick, the Nexus enforces global conservation of DI-bits and
total spin angular momentum.  This conservation applies simultaneously
to all Grid Points, regardless of the spatial separation of the two
aggregates.

The Nexus is \emph{not} a signal traveling between the particles.  It
is a global constraint enforced outside spacetime, at each tick,
simultaneously.  The joint state remains~\eqref{eq:singlet} until one
of the aggregates undergoes a measurement interaction (decoherence
with the DP Sea --- companion CPP-2040d~\citep{cpp2040d}).

\begin{remark}[Why the Nexus is not a local hidden variable]
The Nexus enforces a \emph{global constraint} on the joint DI-bit
state.  It is not a hidden variable carried by individual particles.
Bell's theorem applies to theories where the correlating mechanism
is local (each particle carries its own $\lambda$).  The Nexus is
explicitly global --- it constrains the total DI-bit count over all
Grid Points simultaneously.  Therefore Theorem~\ref{thm:bell} does
not apply to CPP.
\end{remark}

%===========================================================================
\section{Derivation of the Singlet Correlation Function}
\label{sec:correlation}
%===========================================================================

\subsection{Measurement in CPP}

A spin measurement along axis $\hat{a}$ couples the CP aggregate's
DI-bit state to a measurement apparatus (itself a CP aggregate with
a preferred SSV direction $\hat{a}$) via the bit-sea decoherence
process of companion CPP-2040d.  The outcome is:
\begin{equation}
\text{outcome} = \begin{cases}
+1 & \text{with probability } P_+ = |\langle\hat{a}|\psi\rangle|^2, \\
-1 & \text{with probability } P_- = |\langle{-\hat{a}}|\psi\rangle|^2,
\end{cases}
\label{eq:born}
\end{equation}
where $|\hat{a}\rangle$ is the DI-bit eigenstate for the $+1$ outcome
along $\hat{a}$.  This is the Born rule, derived from the DI-bit density
interpretation $P \propto \rho_{\rm bit} = |\psi|^2$ in
companion~C3~\citep{abshier2026born}.

\textbf{The key point}: equation~\eqref{eq:born} uses
$P \propto |\text{projection}|^2$ (Born rule), not
$P \propto |\text{projection}|$ (classical).
This is the mathematical reason CPP escapes the classical bound.

\subsection{Joint probabilities for the singlet}

For the singlet state~\eqref{eq:singlet}, if A is measured along
$\hat{a}$ and B along $\hat{b}$, the Born rule gives:
\begin{align}
P(A{=}+1, B{=}+1) &= \tfrac{1}{2}\sin^2(\theta/2)
  = \tfrac{1-\cos\theta}{4}, \label{eq:ppp}\\
P(A{=}+1, B{=}-1) &= \tfrac{1}{2}\cos^2(\theta/2)
  = \tfrac{1+\cos\theta}{4}, \label{eq:ppm}\\
P(A{=}-1, B{=}+1) &= \tfrac{1}{2}\cos^2(\theta/2)
  = \tfrac{1+\cos\theta}{4}, \label{eq:pmp}\\
P(A{=}-1, B{=}-1) &= \tfrac{1}{2}\sin^2(\theta/2)
  = \tfrac{1-\cos\theta}{4}, \label{eq:pmm}
\end{align}
where $\theta = \angle(\hat{a},\hat{b})$ and we used
$\cos^2(\theta/2) = (1+\cos\theta)/2$, $\sin^2(\theta/2) = (1-\cos\theta)/2$.

\begin{proposition}[Singlet correlation function]
\label{prop:corr}
The expectation value of the product of outcomes is:
\begin{equation}
E(\hat{a},\hat{b}) \;=\; \sum_{a,b \in \{+1,-1\}} ab\,P(A{=}a,B{=}b)
\;=\; -\cos\theta.
\label{eq:corr}
\end{equation}
\end{proposition}

\begin{proof}
\begin{align}
E &= (+1)(+1)\tfrac{1-\cos\theta}{4}
   + (+1)(-1)\tfrac{1+\cos\theta}{4}
   + (-1)(+1)\tfrac{1+\cos\theta}{4}
   + (-1)(-1)\tfrac{1-\cos\theta}{4} \notag\\
  &= \tfrac{(1-\cos\theta) - (1+\cos\theta) - (1+\cos\theta) + (1-\cos\theta)}{4}
  = \tfrac{-4\cos\theta}{4}
  = -\cos\theta. \qed
\end{align}
\end{proof}

%===========================================================================
\section{CHSH = $2\sqrt{2}$}
\label{sec:chsh}
%===========================================================================

\begin{theorem}[Tsirelson bound in CPP]
\label{thm:tsirelson}
With measurement settings
$\hat{a} = 0°$, $\hat{a}' = 90°$, $\hat{b} = 45°$, $\hat{b}' = -45°$,
the CHSH correlator for the CPP singlet state achieves the quantum maximum:
\begin{equation}
S = E(\hat{a},\hat{b}) + E(\hat{a},\hat{b}')
  + E(\hat{a}',\hat{b}) - E(\hat{a}',\hat{b}')
  = -2\sqrt{2}, \quad |S| = 2\sqrt{2}.
\label{eq:tsirelson}
\end{equation}
\end{theorem}

\begin{proof}
The angles between measurement axes:
$\angle(\hat{a},\hat{b}) = 45°$,
$\angle(\hat{a},\hat{b}') = 45°$,
$\angle(\hat{a}',\hat{b}) = 45°$,
$\angle(\hat{a}',\hat{b}') = 135°$.
Using Proposition~\ref{prop:corr}:
\begin{align}
E(\hat{a},\hat{b}) &= -\cos 45° = -1/\sqrt{2}, \notag\\
E(\hat{a},\hat{b}') &= -\cos 45° = -1/\sqrt{2}, \notag\\
E(\hat{a}',\hat{b}) &= -\cos 45° = -1/\sqrt{2}, \notag\\
E(\hat{a}',\hat{b}') &= -\cos 135° = +1/\sqrt{2}. \notag
\end{align}
Therefore:
$S = -1/\sqrt{2} + (-1/\sqrt{2}) + (-1/\sqrt{2}) - (+1/\sqrt{2})
   = -4/\sqrt{2} = -2\sqrt{2}$.
Since $|{-2\sqrt{2}}| = 2\sqrt{2} > 2$, the classical Bell bound is violated.
\end{proof}

%===========================================================================
\section{The No-Signaling Theorem in CPP}
\label{sec:nosig}
%===========================================================================

\begin{theorem}[No-signaling]
\label{thm:nosig}
Alice's marginal measurement probability is independent of Bob's measurement
setting, and vice versa.  No information can be transmitted between the
parties by choice of measurement axis.
\end{theorem}

\begin{proof}
Alice's marginal probability of getting $+1$ (summing over Bob's outcomes):
\begin{equation}
P(A=+1) = P(A{=}+1,B{=}+1) + P(A{=}+1,B{=}-1)
= \frac{1-\cos\theta}{4} + \frac{1+\cos\theta}{4}
= \frac{1}{2}.
\label{eq:nosig_a}
\end{equation}
This is $1/2$ regardless of $\theta$, i.e., regardless of Bob's choice
of measurement axis $\hat{b}$.  By the symmetry of~\eqref{eq:ppp}--\eqref{eq:pmm}
the same holds for Bob's marginal: $P(B=+1) = 1/2$ regardless of $\hat{a}$.
Therefore neither party can signal to the other by their choice of axis.
\end{proof}

\begin{remark}[Why no-signaling holds despite non-locality]
The correlation between outcomes is non-local (it exceeds the classical
bound), yet no signal is transmitted.  This is possible because
Alice and Bob cannot \emph{control} which outcome they get ---
each outcome occurs with probability $1/2$ regardless of the axis
chosen.  The non-local correlation is only visible when the two
parties later \emph{compare} their results (a process requiring
classical communication).
\end{remark}

%===========================================================================
\section{Why the Previous CPP Argument Was Wrong}
\label{sec:previous}
%===========================================================================

The previous version argued: ``correlations are pre-established through
shared bit pool, with local measurements only revealing pre-existing
geometric relationships.''

This is an LHV model.  The ``shared bit pool with pre-existing
geometric relationships'' plays the role of Bell's hidden variable $\lambda$.
The ``local measurement revealing pre-existing relationships'' is exactly
the structure Bell's theorem rules out.

An LHV model with anti-correlated phases on a sphere gives:
\begin{equation}
E_{\rm LHV}(\hat{a},\hat{b}) = 1 - \frac{2\theta}{\pi},
\label{eq:classicalcorr}
\end{equation}
which satisfies $|\mathrm{CHSH}| \leq 2$.  This is demonstrated by the
original Bell (1964) paper.

The corrected CPP model differs in two ways:
\begin{enumerate}
\item The state is \textbf{non-separable} (cannot be factored into
  independent local states).
\item The probability rule is the \textbf{Born rule} ($P = |\langle\hat{b}
  |\psi\rangle|^2$), not the classical projection ($P \propto
  |\hat{b}\cdot\hat{\lambda}|$).
\end{enumerate}
Both differences are required.  The Born rule applied to a separable
state still gives the classical result; the non-separable state with
a classical probability rule also gives the classical result.  Only
the combination gives $E = -\cos\theta$ and $|\mathrm{CHSH}| = 2\sqrt{2}$.

%===========================================================================
\section{Comparison with Other Interpretations of Quantum Mechanics}
\label{sec:compare}
%===========================================================================

\begin{table}[htbp]
\centering
\caption{Comparison of approaches to Bell inequality violation.}
\label{tab:compare}
\begin{tabular}{@{}p{3.2cm}p{3.5cm}p{3.5cm}p{3.5cm}@{}}
\toprule
 & Copenhagen & Many-Worlds & \textbf{CPP (Nexus)} \\
\midrule
Non-local resource & Wavefunction collapse & Global branch structure & Nexus (atemporal global constraint) \\
Mechanism & Observer-induced collapse & Branching at measurement & Bit-sea decoherence (CPP-2040d) \\
Non-separability & Postulated & Postulated & Nexus-maintained DI-bit state \\
Born rule & Postulated & Derived from branch counting & Derived from DI-bit density (C3) \\
Signaling & None & None & None (Theorem~\ref{thm:nosig}) \\
Bell violation & Yes & Yes & Yes (Theorem~\ref{thm:tsirelson}) \\
New prediction & None & None & Lattice corrections at $\nu \sim \nu_P$ \\
\bottomrule
\end{tabular}
\end{table}

CPP is most closely related to objective collapse theories (GRW-style)
in that the collapse mechanism (bit-sea decoherence) is physically
objective.  But it differs fundamentally: GRW introduces a new collapse
postulate; CPP derives collapse from the DP Sea thermalization process
that is already present in the framework for independent reasons
(companion CPP-2040d).

CPP differs from Bohmian mechanics in that there is no pilot wave
guiding each particle's trajectory.  Instead, the Nexus is the
non-local mechanism that maintains the joint state.  Bohmian mechanics
also violates Bell's locality condition (through the pilot wave); CPP
violates it through the Nexus.

%===========================================================================
\section{Falsifiable Predictions}
\label{sec:predictions}
%===========================================================================

\subsection{Standard regime: identical to quantum mechanics}

For all currently achievable experimental frequencies and energies,
CPP predicts correlations identical to standard quantum mechanics:
$E(\hat{a},\hat{b}) = -\cos\theta$.  All loophole-free Bell tests
\citep{hensen2015,giustina2015,shalm2015} are consistent with CPP.

\subsection{Lattice corrections near the Planck scale}

At frequencies approaching the Planck frequency
$\nu_P = 1/t_P \approx 1.855 \times 10^{43}$~Hz, the discrete
600-cell lattice structure introduces corrections to the correlation
function.  The DI-bit phases are defined on a discrete lattice with
spacing $\ell_P$; the phase resolution at frequency $\nu$ is limited
by the lattice to $\delta\phi \sim (\nu/\nu_P)$.  This modifies the
correlation to:
\begin{equation}
E_{\rm CPP}(\hat{a},\hat{b},\nu) \;=\;
-\cos\theta \times \left[1 - \left(\frac{\nu}{\nu_P}\right)^2\right]
+ O\!\left(\frac{\nu}{\nu_P}\right)^4,
\label{eq:correction}
\end{equation}
where the leading correction is $O((\nu/\nu_P)^2)$.

For accessible frequencies the correction is negligibly small:
\begin{equation}
\left(\frac{10^{12}\,\text{Hz}}{1.855\times 10^{43}\,\text{Hz}}\right)^2
\approx 2.9 \times 10^{-63}.
\end{equation}
This is unmeasurable with any foreseeable technology.  The prediction
is genuinely falsifiable in principle --- any theory that predicts a
perfect $|\mathrm{CHSH}| = 2\sqrt{2}$ at all energies (including those
above the Planck scale) is inconsistent with a discrete lattice theory
like CPP --- but it does not provide a near-term experimental target.

\subsection{Entanglement fragility near black hole horizons}

The Nexus maintains the non-separable joint DI-bit state across all
Grid Points.  Near a black hole horizon, the PSR contraction
(companion~C8~\citep{abshier2026gr8}) affects the effective Planck
length and thus the lattice phase resolution.  CPP predicts that
the effective lattice correction becomes:
\begin{equation}
\left(\frac{\nu}{\nu_P}\right)^2 \;\to\;
\left(\frac{\nu}{\nu_P(r)}\right)^2, \quad
\nu_P(r) = \nu_P \times \left(1 + \frac{GM}{rc^2}\right)^{-1},
\end{equation}
so that correlations degrade faster for experiments conducted near
a gravitational source.  This coupling between gravity and entanglement
is a uniquely CPP prediction.

%===========================================================================
\section{Connection to Earlier CPP Papers}
\label{sec:connections}
%===========================================================================

The Bell analysis in this paper depends on three earlier CPP results:

\begin{enumerate}
\item \textbf{Born rule from DI-bit density (companion C3
  \citep{abshier2026born})}: The probability rule
  $P \propto |\psi|^2 = \rho_{\rm bit}$ is what distinguishes the
  quantum correlation $-\cos\theta$ from the classical
  $1 - 2\theta/\pi$.  Without the Born rule derivation, the CPP
  Bell analysis reduces to an LHV model.

\item \textbf{Schr\"{o}dinger equation from DI-bit diffusion
  (CPP-2040b \citep{cpp2040b})}: The linear evolution of the joint
  DI-bit state~\eqref{eq:singlet} between measurements is governed
  by the Schr\"{o}dinger equation.  Non-separability is preserved
  under linear evolution.

\item \textbf{Decoherence from bit-sea thermalization (CPP-2040d
  \citep{cpp2040d})}: The measurement process that selects a definite
  outcome is the irreversible coupling of the CP aggregate to the DP
  Sea.  The specific outcome probabilities follow from the Born rule
  applied to the state just before thermalization.
\end{enumerate}

%===========================================================================
\section{Conclusion}
\label{sec:conclusion}
%===========================================================================

The CPP derivation of Bell inequality violation rests on two pillars:

\begin{enumerate}
\item \textbf{Non-separability}: The singlet DI-bit state cannot be
  factored into independent states for the two particles
  (Proposition~\ref{prop:nonsep}).  The Nexus maintains this
  non-separability as particles separate.  Bell's theorem does not
  apply to non-separable states.

\item \textbf{Born rule}: The probability rule $P = |\langle\hat{b}
  |\psi\rangle|^2$ (derived from DI-bit density in companion~C3)
  gives $E = -\cos\theta$, not $1 - 2\theta/\pi$.  The quantum
  correlation exceeds the classical bound; the classical correlation
  does not.
\end{enumerate}

Together these give $|\mathrm{CHSH}| = 2\sqrt{2}$ at optimal settings
(Theorem~\ref{thm:tsirelson}) and no-signaling (Theorem~\ref{thm:nosig}),
consistent with all loophole-free experiments.

The previous version of this paper used an LHV argument (pre-established
shared bit pool) that Bell's theorem explicitly rules out.  That argument
gives $|\mathrm{CHSH}| \leq 2$ and is incorrect.  The corrected version
identifies the Nexus-maintained non-separable state as the correct
non-local resource, and the Born rule as the correct probability rule.

\textbf{One new prediction} is made: the correlation function acquires
lattice corrections of order $(\nu/\nu_P)^2$ near the Planck frequency,
representing the first point where CPP and standard quantum mechanics
diverge in the Bell context.

\bibliographystyle{unsrt}
\begin{thebibliography}{99}

\bibitem{bell1964}
J.~S. Bell,
``On the Einstein Podolsky Rosen paradox,''
\textit{Physics Physique Fizika} \textbf{1}(3), 195--200 (1964).

\bibitem{clauser1969}
J.~F. Clauser, M.~A. Horne, A.~Shimony, and R.~A. Holt,
``Proposed experiment to test local hidden-variable theories,''
\textit{Phys.\ Rev.\ Lett.} \textbf{23}, 880 (1969).

\bibitem{cirelson1980}
B.~S. Cirel'son,
``Quantum generalizations of Bell's inequality,''
\textit{Lett.\ Math.\ Phys.} \textbf{4}, 93--100 (1980).

\bibitem{hensen2015}
B.~Hensen \textit{et al.},
``Loophole-free Bell inequality violation using electron spins
separated by 1.3 kilometres,''
\textit{Nature} \textbf{526}, 682--686 (2015).

\bibitem{giustina2015}
M.~Giustina \textit{et al.},
``Significant-loophole-free test of Bell's theorem with entangled photons,''
\textit{Phys.\ Rev.\ Lett.} \textbf{115}, 250401 (2015).

\bibitem{shalm2015}
L.~K. Shalm \textit{et al.},
``Strong loophole-free test of local realism,''
\textit{Phys.\ Rev.\ Lett.} \textbf{115}, 250402 (2015).

\bibitem{abshier2026born}
T.~L. Abshier,
``The Born Rule from DI-bit Density in Conscious Point Physics,''
companion~C3 (2026).

\bibitem{abshier2026zbw}
T.~L. Abshier,
``Inertial Mass from Zitterbewegung,''
companion~C4 (2026).

\bibitem{abshier2026gr8}
T.~L. Abshier,
``Strong-Field General Relativity and the CPP Field Equation,''
companion~C8 (2026).

\bibitem{cpp2040b}
T.~L. Abshier and Grok,
``Quantum Mechanics in CPP: Schr\"{o}dinger Equation from
DI-bit Diffusion,'' CPP-2040b (2026).

\bibitem{cpp2040d}
T.~L. Abshier and Grok,
``Quantum Mechanics in CPP: Measurement Problem from
Bit-Sea Thermalization,'' CPP-2040d (2026).

\bibitem{einstein1935}
A.~Einstein, B.~Podolsky, and N.~Rosen,
``Can quantum-mechanical description of physical reality
be considered complete?''
\textit{Phys.\ Rev.} \textbf{47}, 777 (1935).

\end{thebibliography}

%===========================================================================
\appendix
%===========================================================================

\section{Why the Classical Anti-Correlation Model Fails}
\label{app:classical}

For completeness, here is the explicit calculation showing why the
``pre-established anti-correlated phases'' (LHV) model gives the
wrong result.

In the classical model, each particle carries a unit vector $\hat{\lambda}$
on the sphere.  Particle~A carries $+\hat{\lambda}$; particle~B carries
$-\hat{\lambda}$.  The outcomes are:
$A = \text{sign}(\hat{a}\cdot\hat{\lambda})$,
$B = -\text{sign}(\hat{b}\cdot\hat{\lambda})$.

The correlation is:
\begin{equation}
E_{\rm cl}(\hat{a},\hat{b})
= \int_{\mathcal{S}^2} \text{sign}(\hat{a}\cdot\hat{\lambda})\,
  (-\text{sign}(\hat{b}\cdot\hat{\lambda}))\,\frac{d^2\lambda}{4\pi}
= 1 - \frac{2\theta}{\pi},
\label{eq:classicalcalc}
\end{equation}
where $\theta = \angle(\hat{a},\hat{b})$.  This is derived by noting
that the product of signs is $-1$ when $\hat{\lambda}$ falls in the
region between the two half-spaces defined by $\hat{a}$ and $\hat{b}$,
which has fractional area $\theta/\pi$.

With this model at the optimal CHSH settings ($\theta_{ab} = \theta_{ab'}
= \theta_{a'b} = 45°$, $\theta_{a'b'} = 135°$):
\begin{align}
E_{\rm cl}(\hat{a},\hat{b}) &= 1 - 2(45°)/180° = 1/2, \\
E_{\rm cl}(\hat{a},\hat{b}') &= 1/2, \\
E_{\rm cl}(\hat{a}',\hat{b}) &= 1/2, \\
E_{\rm cl}(\hat{a}',\hat{b}') &= 1 - 2(135°)/180° = -1/2.
\end{align}
\begin{equation}
|S_{\rm cl}| = |1/2 + 1/2 + 1/2 - (-1/2)| = 2.
\end{equation}
The classical model saturates the Bell bound exactly.  It never exceeds 2.
The quantum prediction gives $2\sqrt{2} \approx 2.828 > 2$.

This demonstrates explicitly that adding ``shared phases'' or
``pre-established correlations'' to a deterministic local model cannot
produce quantum correlations.  The quantum value requires
\emph{non-separability} and the \emph{Born rule} --- neither of which
is present in the classical model.


\end{document}
